\documentclass[journal]{IEEEtran}

% Packages
\usepackage[utf8]{inputenc}
\usepackage{amsmath,amssymb}
\usepackage{graphicx}
\usepackage{cite}

\begin{document}

\title{Modular versus Monolithic Neural Architectures for Gene Regulatory
Network Inference: Capacity, Robustness, and Representational Efficiency}

\author{Evint Leovonzko%
\thanks{E. Leovonzko is with [Department], [University], [City, Country].
E-mail: [email]. Manuscript received 2026.}}

\maketitle

\begin{abstract}
Gene regulatory network (GRN) inference --- determining which transcription
factors regulate which target genes from single-cell RNA-sequencing data --- is
a core problem in computational biology. A fundamental architectural choice is
whether to encode transcription factor (TF) and gene representations
independently (two-tower/dual-encoder) or jointly (cross-encoder). We present
a parameter-matched comparison of both architectures on a human brain
single-nucleus RNA-seq dataset (39K TF--gene pairs, 5.58M parameters each)
under balanced (1:1) and realistic (5:1) negative-to-positive sampling ratios.
The cross-encoder achieves an AUROC of 0.9040 versus 0.8097 for the two-tower
under balanced training --- a gap of 9.4 points consistent across five random
seeds. Under 5:1 negative sampling, the two-tower AUROC degrades by 6.6 points
(0.8097~$\to$~0.7434) while the cross-encoder remains stable (0.9040~$\to$~0.9150),
demonstrating that cosine similarity scoring is geometrically sensitive to
class imbalance. A structured neuron pruning analysis on the trained two-tower
model shows that post-hoc AUROC retention never falls below 99.96\% of
baseline at any sparsity level tested up to 90\% (51 of 512 neurons retained
per tower), revealing that the two-tower fails to utilize its nominal
parameter budget. This redundancy provides a mechanistic account of the
cross-encoder's advantage and suggests that a 90\%-sparse two-tower matches
baseline performance without any fine-tuning.
\end{abstract}

\begin{IEEEkeywords}
gene regulatory network inference, two-tower architecture, cross-encoder,
dual-encoder, neural network pruning, single-cell RNA-seq, representation
learning, link prediction, class imbalance
\end{IEEEkeywords}

% ── I. INTRODUCTION ──────────────────────────────────────────────────────────
\section{Introduction}
\label{sec:intro}

Gene regulatory network (GRN) inference --- the problem of identifying which
transcription factors (TFs) regulate which target genes --- is a fundamental
challenge in computational biology. Understanding transcriptional regulatory
relationships is essential for deciphering developmental programs, disease
mechanisms, and therapeutic targets. With the availability of large-scale
single-nucleus RNA-sequencing (snRNA-seq) data, machine learning methods have
emerged that cast GRN inference as binary link prediction: given a candidate
TF--gene pair and their expression profiles, predict whether a regulatory
relationship exists~\cite{aibar2017,huynh2010}.

A central architectural question in link prediction is whether entity
representations should be computed independently (a \emph{two-tower} or
dual-encoder model~\cite{huang2013,bromley1993}) or jointly (a
\emph{cross-encoder}~\cite{nogueira2019,humeau2020}). Two-tower models score
interactions via cosine similarity of independently computed encodings,
enabling efficient inference at scale. Cross-encoders process pairs jointly,
allowing arbitrary interaction features to be learned before scoring. In
information retrieval, cross-encoders consistently outperform two-tower models
when precise ranking quality is needed~\cite{nogueira2019,humeau2020}, but
whether this advantage holds in computational biology --- where entities have
distinct biological identities and the regulatory signal is low-dimensional ---
has not been established under controlled, parameter-matched conditions.

Three open questions motivate this work: (1)~Does a cross-encoder achieve
higher discriminative AUROC than a parameter-matched two-tower for TF--gene
link prediction? (2)~Is the two-tower's cosine similarity scoring
disproportionately sensitive to class imbalance relative to a cross-encoder?
(3)~Does the trained two-tower actually utilize its nominal hidden-layer
capacity, or is much of it redundant?

This paper makes three contributions:
\begin{itemize}
\item \textbf{C1:} A parameter-matched comparison (5.58M parameters each)
  of two-tower and cross-encoder MLPs on human brain snRNA-seq data under
  both balanced (1:1) and realistic (5:1) negative-to-positive sampling ratios.
\item \textbf{C2:} Demonstration that the two-tower degrades by 6.6~AUROC
  points under 5:1 negative sampling --- approximating real-world regulatory
  database sparsity --- while the cross-encoder remains stable.
\item \textbf{C3:} A structured neuron pruning analysis showing that
  post-hoc AUROC retention is $\geq$0.9996 at all sparsity levels up to
  90\%, revealing that the two-tower underutilizes its parameter budget
  due to the geometric constraints of cosine similarity scoring.
\end{itemize}

The remainder of this paper is organized as follows.
Section~\ref{sec:related} reviews related work.
Section~\ref{sec:methods} describes the dataset, architectures, training
protocol, and pruning procedure.
Section~\ref{sec:results} presents experimental results.
Section~\ref{sec:discussion} discusses implications and limitations.
Section~\ref{sec:conclusion} concludes.

\end{document}
